\subsection{Reflexive spaces}

DEFINITION 3. --- A locally convex space E is said to be reflexive if the canonical mapping \(c_{\mathrm{E}}\) from E into \(\mathbf{E}''\) is a topological vector space isomorphism from E onto the strong dual of \(\mathbf{E}_b'\) .

In particular, a reflexive space is semi-reflexive, hence Hausdorff.

PROPOSITION 4. --- The strong dual of a reflexive space is reflexive.

This follows immediately from def. 3.

THEOREM 2. --- In order that a locally convex Hausdorff space E be reflexive, it is necessary and sufficient that it is barrelled and that every bounded subset of E is relatively compact for the weakened topology \(\sigma(E, E')\) .

By th. 1 (IV, p.~15), this is the same as saying that E is reflexive if and only if it is semi-reflexive and barrelled.

If E is reflexive, \(E_{b}^{\prime}\) is reflexive (prop. 4) and consequently E is barrelled (IV, p.~15, th. 1). Conversely, if E is semi-reflexive and barrelled, \(c_{E}\) is a bijection and is bicontinuous by IV, p.~15, prop. 2, hence E is reflexive.

Remarks. --- * 1) Let E be an infinite dimensional real hilbertian space. Let F denote the space E endowed with the weakened topology. The spaces E and F have the same dual E', and E is a reflexive Banach space (V, p.~17). Consequently, F is semi-reflexive.

However, the strong topology and the weakened topology on E are distinct, hence F is not reflexive. *

\begin{enumerate}
\def\labelenumi{\arabic{enumi})}
\setcounter{enumi}{1}
\tightlist
\item
  Let E be a reflexive space and M a closed vector subspace of E. It may happen that neither M nor E/M are reflexive spaces (IV, p.~63, exerc. 10). * For the case of normed spaces, see prop. 7 of IV, p.~17. *
\end{enumerate}

